\GXNSFBodyText

\noindent\textbf{实施地点（占位示例）}\par
项目拟在依托单位的计算实验环境和合规授权的在线数据环境中实施，主要工作包括公开资料核验、事件标注、模型训练与结果复核，不开展针对佐佐木希本人或普通平台用户的线下追踪和干预实验。正式申报时应替换为真实实验室名称、地址、设备使用地点及数据访问环境。

\noindent\textbf{实施规模}\par
研究规模按“一个贯通案例、若干同期同构人物对照、一个区域组件测试场景”组织。佐佐木希公开资料用于贯通全流程，同构人物案例用于完整模型的外部效度检验；广西—东盟样本以单一文化内容 IP 或单一区域品牌的固定观察窗为单位，只测试多语言实体对齐、平台校准和扰动窗口组件。具体平台数、时间跨度、事件条数和人工标注量在预实验后依据统计效能、数据许可和工作量确定；模板不预填未经证实的“十万条舆情”等数字。
